\chapter{Nền tảng lý thuyết và công nghệ}

\section{Micro-learning và Bloom Taxonomy}

Micro-learning là cách tổ chức nội dung học tập thành các đơn vị nhỏ, tập
trung vào một khái niệm hoặc kỹ năng cụ thể. Trong bối cảnh LMS, mỗi đơn vị có
thể được biểu diễn dưới dạng lesson card, flashcard, câu hỏi tự kiểm tra hoặc
một đoạn giải thích ngắn. Cách tiếp cận này phù hợp với nhu cầu ôn tập nhanh,
học lặp lại và theo dõi tiến độ theo từng phần nhỏ của học phần.

Đối với đề tài, Micro-learning không chỉ là định dạng hiển thị mà còn ảnh
hưởng đến pipeline xử lý tài liệu. Tài liệu dài cần được chia thành các chunk
có ngữ cảnh đủ rõ, sau đó chuyển thành nội dung học tập ngắn có liên kết nguồn
và có thể kiểm duyệt bởi giảng viên.

Bloom Taxonomy được sử dụng để phân loại mức độ nhận thức của Learning
Outcome và câu hỏi. Sáu mức gồm Remember, Understand, Apply, Analyze,
Evaluate và Create. Với quiz trắc nghiệm trong phạm vi hệ thống, các mức
Remember, Understand, Apply và Analyze là phù hợp nhất; Evaluate và Create
thường cần dạng tự luận hoặc bài tập dự án.

Hệ thống sử dụng Bloom level như metadata và ràng buộc sinh câu hỏi. Ví dụ,
Learning Outcome có động từ ``trình bày'', ``liệt kê'' thường gắn với
Remember; ``giải thích'', ``phân biệt'' thường gắn với Understand; ``áp
dụng'', ``tính toán'' thường gắn với Apply; ``phân tích'', ``so sánh'' thường
gắn với Analyze hoặc Evaluate tùy ngữ cảnh. Bloom level ảnh hưởng đến prompt,
độ khó câu hỏi và rubric đánh giá chất lượng quiz.

\section{Mô hình ngôn ngữ lớn (LLM) và Generative AI}

Mô hình ngôn ngữ lớn (Large Language Model --- LLM) là mô hình học sâu có khả
năng hiểu, sinh và biến đổi ngôn ngữ tự nhiên. Trong hệ thống, LLM được dùng
cho các tác vụ cần diễn giải tri thức:
\begin{itemize}
    \item Tóm tắt chunk thành Micro-Content hoặc Flashcard.
    \item Sinh câu hỏi trắc nghiệm, đáp án đúng, distractor và giải thích.
    \item Chuẩn hóa hoặc trích xuất thông tin từ đề cương môn học khi parser
          heuristic không đủ.
    \item Hỗ trợ RAG chatbot hoặc trợ lý học tập trong hướng mở rộng.
\end{itemize}

Để giảm phụ thuộc vào một provider cụ thể, hệ thống đặt LLM sau interface
\texttt{LLMClient}. Use case chỉ yêu cầu một contract: nhận prompt có cấu trúc
và trả về dữ liệu đúng schema. Provider như Gemini, GPT hoặc mô hình mã nguồn
mở có thể được thay đổi ở tầng adapter mà không làm thay đổi nghiệp vụ cốt
lõi.

Prompt được thiết kế theo nguyên tắc: nêu vai trò, nêu nhiệm vụ, đưa dữ liệu
nguồn, yêu cầu định dạng đầu ra và ràng buộc chất lượng. Với quiz generation,
prompt cần chứa context chunk, Learning Outcome nếu có, Bloom level mục tiêu,
số lượng câu hỏi, ngôn ngữ và schema JSON đầu ra gồm question, options,
correct answer, explanation, source chunk và rationale.

Sau khi LLM trả về, hệ thống cần validate JSON, kiểm tra số lượng option, đáp
án đúng, độ dài câu hỏi và giải thích. Nếu JSON lỗi, adapter có thể retry với
prompt sửa lỗi. Nếu nội dung vẫn không đạt, kết quả được đánh dấu cần review
thay vì publish tự động.

\section{Kiến trúc RAG và GraphRAG trong giáo dục}

Retrieval-Augmented Generation (RAG) kết hợp truy xuất thông tin với LLM.
Thay vì để LLM trả lời dựa hoàn toàn trên tri thức nội tại, hệ thống cung cấp
context từ tài liệu đã xử lý. Cách tiếp cận này giúp tăng khả năng bám sát
tài liệu gốc, cho phép liên kết nguồn chunk và giảm nguy cơ sinh nội dung
không thuộc học phần.

Pipeline RAG trong hệ thống gồm các bước cơ bản:
\begin{enumerate}
    \item Chuẩn hóa query hoặc yêu cầu sinh nội dung.
    \item Tạo embedding cho query.
    \item Tìm top-K chunk liên quan trong Qdrant.
    \item Lọc theo course, document, chapter hoặc quyền truy cập LMS.
    \item Lắp prompt từ instruction, Learning Outcome, Bloom level và context
          chunk.
    \item Gọi LLM để sinh câu trả lời, card hoặc quiz.
\end{enumerate}

GraphRAG mở rộng RAG bằng cách dùng Knowledge Graph để bổ sung quan hệ cấu
trúc. Trong giáo dục, tri thức không chỉ là đoạn văn bản mà còn có quan hệ
giữa Course, Chapter, Learning Outcome, Assessment và nội dung học tập. Khi
giảng viên yêu cầu sinh quiz cho một LO cụ thể, hệ thống không chỉ tìm chunk
gần query bằng vector similarity mà còn ưu tiên chunk có quan hệ hỗ trợ LO đó.

Đề tài tiếp cận GraphRAG theo hướng curriculum-aware. Knowledge Graph được
dùng để lọc, tăng trọng số hoặc giải thích mối liên hệ giữa nội dung và chuẩn
đầu ra. Graph reasoning nhiều bước là hướng mở rộng, còn trọng tâm luận văn là
xây dựng nền tảng truy xuất có ý thức về chương trình học.

\section{Chuẩn tích hợp hệ thống học tập (LMS và LTI 1.3)}

Learning Management System (LMS) như Open edX hoặc Moodle đóng vai trò nền
tảng quản lý khóa học, người dùng, tài liệu và hoạt động học tập. Thay vì xây
dựng một LMS mới, đề tài thiết kế hệ thống như một công cụ học tập bên ngoài
có thể nhúng vào LMS hiện hữu.

Learning Tools Interoperability (LTI) 1.3 là chuẩn cho phép LMS (platform)
khởi chạy một công cụ bên ngoài (tool) với thông tin người dùng, course
context và quyền truy cập đã được xác thực. Luồng LTI 1.3 dựa trên OIDC login
và JWT đã ký. Khi người dùng mở activity từ LMS, LMS gửi launch request đến
tool; tool kiểm tra issuer, client ID, deployment ID, nonce, chữ ký và claim
người dùng trước khi tạo session nội bộ.

Trong thiết kế của luận văn, Open edX là nền tảng tích hợp chính. Moodle được
phân tích như nền tảng tương thích vì cùng hỗ trợ LTI external tool. Cách tích
hợp này giúp hệ thống AI có thể được triển khai bên cạnh LMS mà không can
thiệp sâu vào lõi của từng nền tảng.

\section{Nền tảng kiến trúc và công nghệ lõi}

\subsection{Hexagonal Architecture}

Hexagonal Architecture tách nghiệp vụ cốt lõi khỏi chi tiết hạ tầng. Phần
domain/application chỉ phụ thuộc vào các port trừu tượng, còn các adapter cụ
thể chịu trách nhiệm giao tiếp với PostgreSQL, Qdrant, Neo4j, MinIO, Redis,
LLM provider hoặc LMS.

Trong đề tài, cách tách lớp này giúp các use case như ProcessDocument,
SearchChunks, IngestCurriculum hoặc GenerateCurriculumQuiz không bị phụ thuộc
trực tiếp vào công nghệ lưu trữ hay provider AI. Khi cần thay Qdrant bằng
vector database khác hoặc thay LLM provider, hệ thống chỉ cần thay adapter phù
hợp.

\subsection{Vector Database}

Vector database lưu embedding của chunk để phục vụ semantic search và RAG.
Mỗi chunk sau khi được trích xuất từ tài liệu sẽ được chuyển thành vector số,
kèm metadata như \texttt{chunk\_id}, \texttt{document\_id},
\texttt{course\_id}, heading path, ngôn ngữ và vị trí trong tài liệu.

Qdrant được chọn cho giai đoạn luận văn vì có Docker image chính thức, API
đơn giản, hỗ trợ payload filtering và phù hợp với môi trường tự host. Vector
database là thành phần quyết định chất lượng truy xuất nên cần được đánh giá
bằng các chỉ số như Recall@K, MRR và latency.

\subsection{Knowledge Graph}

Knowledge Graph biểu diễn tri thức dưới dạng node và relationship. Với bài
toán giáo dục, các node chính gồm Course, Chapter, LearningOutcome,
Assessment và Chunk. Các quan hệ như \texttt{COURSE\_HAS\_CHAPTER},
\texttt{CHAPTER\_HAS\_LO}, \texttt{ASSESSMENT\_MEASURES\_LO} và
\texttt{CHUNK\_SUPPORTS\_LO} giúp hệ thống hiểu nội dung nào hỗ trợ chuẩn đầu
ra nào.

Neo4j được sử dụng như graph database để lưu projection tri thức. PostgreSQL
vẫn là source of truth cho dữ liệu nghiệp vụ, còn Neo4j tối ưu cho truy vấn
quan hệ phục vụ curriculum-aware retrieval, GraphRAG và adaptive learning về
sau.

\subsection{Event-driven / Async Processing}

Xử lý tài liệu, sinh embedding và gọi LLM là các tác vụ tốn thời gian, không
phù hợp với request/response đồng bộ. Vì vậy hệ thống sử dụng mô hình xử lý
bất đồng bộ thông qua queue và worker. Sau khi giảng viên upload tài liệu,
backend lưu file, tạo metadata và enqueue job; worker thực hiện parse,
chunking, embedding, generation và cập nhật trạng thái pipeline.

Mô hình event-driven/async giúp frontend phản hồi nhanh, pipeline có thể retry
khi lỗi tạm thời và worker có thể scale độc lập theo số lượng tài liệu. Đồng
thời, outbox/event flow cho phép đồng bộ dữ liệu từ PostgreSQL sang Neo4j mà
không làm chậm luồng nghiệp vụ chính.
